pdfoutput=1
% Options for packages loaded elsewhere
\PassOptionsToPackage{unicode}{hyperref}
\PassOptionsToPackage{hyphens}{url}
\pdfoutput=1
\documentclass[
  12pt,
]{article}
\usepackage{xcolor}
\usepackage{amsmath,amssymb}
\setcounter{secnumdepth}{-\maxdimen} % remove section numbering
\usepackage{iftex}
\ifPDFTeX
  \usepackage[T1]{fontenc}
  \usepackage{textcomp} % provide euro and other symbols
\else % if luatex or xetex
  \usepackage{unicode-math} % this also loads fontspec
  \defaultfontfeatures{Scale=MatchLowercase}
  \defaultfontfeatures[\rmfamily]{Ligatures=TeX,Scale=1}
\fi
\usepackage{lmodern}
\ifPDFTeX\else
  % xetex/luatex font selection
\fi
% Use upquote if available, for straight quotes in verbatim environments
\IfFileExists{upquote.sty}{\usepackage{upquote}}{}
\IfFileExists{microtype.sty}{% use microtype if available
  \usepackage[]{microtype}
  \UseMicrotypeSet[protrusion]{basicmath} % disable protrusion for tt fonts
}{}
\makeatletter
\@ifundefined{KOMAClassName}{% if non-KOMA class
  \IfFileExists{parskip.sty}{%
    \usepackage{parskip}
  }{% else
    \setlength{\parindent}{0pt}
    \setlength{\parskip}{6pt plus 2pt minus 1pt}}
}{% if KOMA class
  \KOMAoptions{parskip=half}}
\makeatother
\usepackage{longtable,booktabs,array}
\usepackage{calc} % for calculating minipage widths
% Correct order of tables after \paragraph or \subparagraph
\usepackage{etoolbox}
\makeatletter
\patchcmd\longtable{\par}{\if@noskipsec\mbox{}\fi\par}{}{}
\makeatother
% Allow footnotes in longtable head/foot
\IfFileExists{footnotehyper.sty}{\usepackage{footnotehyper}}{\usepackage{footnote}}
\makesavenoteenv{longtable}
\setlength{\emergencystretch}{3em} % prevent overfull lines
\providecommand{\tightlist}{%
  \setlength{\itemsep}{0pt}\setlength{\parskip}{0pt}}
\usepackage{bookmark}
\IfFileExists{xurl.sty}{\usepackage{xurl}}{} % add URL line breaks if available
\urlstyle{same}
\hypersetup{
  hidelinks,
  pdfcreator={LaTeX via pandoc}}

\author{SCX}
\date{2026}

\begin{document}

\section{Lemma A: Bahadur-Rao Theorem -- Bernoulli Case, Complete
Derivation}\label{lemma-a-bahadur-rao-theorem-bernoulli-case-complete-derivation}

\begin{quote}
\textbf{Target}: Prove the Bahadur-Rao (1960) theorem specialized to
i.i.d. Bernoulli(p) random variables, with explicit constants and error
bounds, and extend to the lower tail and the non-i.i.d. setting.
\end{quote}

\begin{center}\rule{0.5\linewidth}{0.5pt}\end{center}

\subsection{A.1 Problem Setup}\label{a.1-problem-setup}

Let
\(X_1, \ldots, X_M \stackrel{\text{i.i.d.}}{\sim} \text{Bernoulli}(p)\)
with \(p \in (0,1)\). Define the sample mean

\[\bar{X}_M = \frac{1}{M}\sum_{i=1}^M X_i.\]

For a threshold \(\theta > p\), we are interested in the large deviation
probability

\[P_M(\theta) := \mathbb{P}(\bar{X}_M \geq \theta).\]

\begin{center}\rule{0.5\linewidth}{0.5pt}\end{center}

\subsection{A.2 Cumulant Generating Function and Cramer
Transform}\label{a.2-cumulant-generating-function-and-cramer-transform}

The moment generating function of a single Bernoulli(p) variable \(X\)
is

\[M(\lambda) = \mathbb{E}[e^{\lambda X}] = 1-p + p e^\lambda, \quad \lambda \in \mathbb{R},\]

so the cumulant generating function (CGF) is

\[\psi(\lambda) = \log M(\lambda) = \log(1-p + p e^\lambda).\]

The first two derivatives of \(\psi\) are:

\[\psi'(\lambda) = \frac{p e^\lambda}{1-p + p e^\lambda},\]

\[\psi''(\lambda) = \frac{p e^\lambda (1-p + p e^\lambda) - (p e^\lambda)^2}{(1-p + p e^\lambda)^2}
= \frac{p(1-p)e^\lambda}{(1-p + p e^\lambda)^2}.\]

The rate function (Cramer transform) is the Legendre-Fenchel conjugate
of \(\psi\):

\[I(\theta) = \sup_{\lambda \in \mathbb{R}} \{\lambda\theta - \psi(\lambda)\}, \qquad \theta \in [0,1].\]

For \(\theta \in (0,1)\), the supremum is attained at the unique
\(\lambda^*\) solving \(\psi'(\lambda) = \theta\):

\[\frac{p e^{\lambda^*}}{1-p + p e^{\lambda^*}} = \theta.\]

\textbf{Proposition A.1 (Explicit saddlepoint)}. \emph{The equation
\(\psi'(\lambda^*) = \theta\) has the unique solution}

\[\lambda^* = \log\frac{\theta(1-p)}{p(1-\theta)}.\]

\emph{Proof.} From \(\psi'(\lambda^*) = \theta\):

\[
\frac{p e^{\lambda^*}}{1-p + p e^{\lambda^*}} = \theta
\;\Longrightarrow\; p e^{\lambda^*} = \theta(1-p + p e^{\lambda^*})
\;\Longrightarrow\; p e^{\lambda^*}(1-\theta) = \theta(1-p)
\;\Longrightarrow\; e^{\lambda^*} = \frac{\theta(1-p)}{p(1-\theta)}.
\]

Taking logs gives the result. \(\square\)

\textbf{Proposition A.2 (Rate function = KL divergence)}. \emph{For
\(\theta \in [0,1]\),}

\[I(\theta) = \theta\log\frac{\theta}{p} + (1-\theta)\log\frac{1-\theta}{1-p} = \text{KL}(\theta \| p).\]

\emph{Proof.} Substituting \(\lambda^*\) into
\(\lambda\theta - \psi(\lambda)\):

\[
\begin{aligned}
I(\theta) &= \lambda^*\theta - \log(1-p + p e^{\lambda^*}) \\
&= \theta\log\frac{\theta(1-p)}{p(1-\theta)} - \log\!\left(1-p + p\cdot\frac{\theta(1-p)}{p(1-\theta)}\right) \\
&= \theta\log\frac{\theta(1-p)}{p(1-\theta)} - \log\!\left(\frac{(1-p)(1-\theta) + \theta(1-p)}{1-\theta}\right) \\
&= \theta\log\frac{\theta(1-p)}{p(1-\theta)} - \log\!\left(\frac{1-p}{1-\theta}\right) \\
&= \theta\log\frac{\theta}{p} + (1-\theta)\log\frac{1-\theta}{1-p} = \text{KL}(\theta \| p). \qquad \square
\end{aligned}
\]

\begin{center}\rule{0.5\linewidth}{0.5pt}\end{center}

\subsection{A.3 Exponential Tilting and the Tilted
Distribution}\label{a.3-exponential-tilting-and-the-tilted-distribution}

Define the exponentially tilted (Esscher-transformed) probability
measure \(\mathbb{P}_{\lambda^*}\) by

\[\frac{d\mathbb{P}_{\lambda^*}}{d\mathbb{P}} = \exp\!\big(\lambda^* S_M - M\psi(\lambda^*)\big),\]

where \(S_M = \sum_{i=1}^M X_i\). Under \(\mathbb{P}_{\lambda^*}\), the
\(X_i\) are i.i.d. with tilted marginal

\[\mathbb{P}_{\lambda^*}(X_i = x) = \frac{e^{\lambda^* x}}{1-p + p e^{\lambda^*}} \cdot \mathbb{P}(X_i = x), \qquad x \in \{0,1\}.\]

\textbf{Proposition A.3 (Tilted distribution is Bernoulli(\(\theta\)))}.
\emph{Under \(\mathbb{P}_{\lambda^*}\),}

\[X_i \sim \text{Bernoulli}(\theta), \quad \text{with } \mathbb{E}_{\lambda^*}[X_i] = \theta,\;
\operatorname{Var}_{\lambda^*}(X_i) = \sigma^2(\theta) = \theta(1-\theta).\]

\emph{Proof.} Using \(e^{\lambda^*} = \theta(1-p)/[p(1-\theta)]\):

\[
\begin{aligned}
\mathbb{P}_{\lambda^*}(X_i=1) &= \frac{e^{\lambda^*} p}{1-p + p e^{\lambda^*}}
= \frac{\frac{\theta(1-p)}{p(1-\theta)}\cdot p}{\frac{1-p}{1-\theta}}
= \frac{\theta}{1-\theta} \cdot (1-\theta) = \theta.
\end{aligned}
\]

Thus \(\mathbb{P}_{\lambda^*}(X_i=1)=\theta\), so
\(X_i \sim \text{Bern}(\theta)\). The variance is \(\theta(1-\theta)\).

Alternatively, from the CGF:

\[\psi''(\lambda^*) = \frac{p(1-p)e^{\lambda^*}}{(1-p + p e^{\lambda^*})^2}.\]

Substituting \(e^{\lambda^*}\) and
\(1-p + p e^{\lambda^*} = (1-p)/(1-\theta)\):

\[
\psi''(\lambda^*) = \frac{p(1-p)\cdot\frac{\theta(1-p)}{p(1-\theta)}}{\left(\frac{1-p}{1-\theta}\right)^2}
= \frac{\theta(1-\theta)(1-p)^2/(1-\theta)^2}{(1-p)^2/(1-\theta)^2}
= \theta(1-\theta). \qquad \square
\]

\begin{center}\rule{0.5\linewidth}{0.5pt}\end{center}

\subsection{A.4 The Bahadur-Rao Theorem for
Bernoulli}\label{a.4-the-bahadur-rao-theorem-for-bernoulli}

\subsubsection{A.4.1 Statement}\label{a.4.1-statement}

\textbf{Theorem A.4 (Bahadur-Rao, 1960; Bernoulli case)}. \emph{Let
\(X_1,\ldots,X_M \stackrel{i.i.d.}{\sim} \text{Bernoulli}(p)\) and fix
\(\theta > p\). Then}

\[
\mathbb{P}(\bar{X}_M \geq \theta) = \frac{\exp(-M \cdot \text{KL}(\theta \| p))}
{\lambda^* \sqrt{2\pi M \cdot \theta(1-\theta)}}
\Bigl(1 + \varepsilon_M\Bigr),
\]

\emph{where \(\lambda^* = \log\frac{\theta(1-p)}{p(1-\theta)} > 0\) and}

\[|\varepsilon_M| \leq \frac{C(p,\theta)}{M}\]

\emph{for an explicit constant \(C(p,\theta)\) given in Proposition A.5
below.}

\textbf{Remark}. The result also holds for the strict inequality
\(\bar{X}_M > \theta\), since \(\mathbb{P}(\bar{X}_M = \theta) = 0\) for
non-integer \(M\theta\), and for integer \(M\theta\) the correction is
absorbed into the \(O(1/M)\) term.

\subsubsection{A.4.2 Proof via Stirling's Formula (Direct Binomial Tail
Expansion)}\label{a.4.2-proof-via-stirlings-formula-direct-binomial-tail-expansion}

Since \(S_M = \sum_{i=1}^M X_i \sim \text{Binomial}(M,p)\), we can work
directly with the binomial distribution. Let
\(k = \lceil M\theta \rceil\). Since \(\theta > p\) and \(\theta\) is
fixed, we have \(k > Mp\) for all sufficiently large \(M\).

The binomial tail probability is

\[P_M(\theta) = \mathbb{P}(S_M \geq k) = \sum_{j=k}^M a_j, \qquad
a_j = \binom{M}{j} p^{\,j}(1-p)^{M-j}.\]

\emph{Step 1: Stirling's formula with explicit remainder.}

For any \(n \in \mathbb{N}\), Robbins' refinement of Stirling gives

\[\sqrt{2\pi n}\left(\frac{n}{e}\right)^n \exp\!\left(\frac{1}{12n+1}\right) < n! < \sqrt{2\pi n}\left(\frac{n}{e}\right)^n \exp\!\left(\frac{1}{12n}\right).\]

Thus \(n! = \sqrt{2\pi n}\,(n/e)^n\,e^{r_n}\) with
\(|r_n| \leq 1/(12n)\) for \(n \geq 1\).

Applying this to \(\binom{M}{k} = M!/(k!\,(M-k)!)\):

\[
\binom{M}{k} = \sqrt{\frac{M}{2\pi k(M-k)}}\,
\frac{M^M}{k^k (M-k)^{M-k}}\,
\exp\!\big(r_M - r_k - r_{M-k}\big).
\]

The error term satisfies

\[|r_M - r_k - r_{M-k}| \leq \frac{1}{12M} + \frac{1}{12k} + \frac{1}{12(M-k)}
\leq \frac{1}{4\min(k, M-k)}.\]

Define \(\theta_M = k/M\). Writing \(H(x) = -x\log x - (1-x)\log(1-x)\)
(binary entropy):

\[\binom{M}{k} = \frac{\exp\!\big(M\cdot H(\theta_M)\big)}{\sqrt{2\pi M\cdot\theta_M(1-\theta_M)}}
\cdot \big(1 + \delta_M^{(1)}\big), \qquad |\delta_M^{(1)}| \leq \frac{C_1}{M}\]

for an explicit constant \(C_1 = C_1(p,\theta)\) that depends on \(p\)
and \(\theta\) but not on \(M\). (Here we bound
\(|e^{r_M-r_k-r_{M-k}} - 1|\) using \(|e^u-1| \leq |u|e^{|u|}\).)

\emph{Step 2: Leading term \(a_k\).}

\[
\begin{aligned}
a_k &= \binom{M}{k} p^k (1-p)^{M-k} \\
&= \frac{\exp\!\big(-M\cdot I(\theta_M)\big)}{\sqrt{2\pi M\cdot\theta_M(1-\theta_M)}}
\cdot \big(1 + \delta_M^{(1)}\big),
\end{aligned}
\]

where
\(I(x) = x\log\frac{x}{p} + (1-x)\log\frac{1-x}{1-p} = \text{KL}(x\|p)\).
Since \(|\theta_M - \theta| \leq 1/M\), a Taylor expansion gives

\[I(\theta_M) = I(\theta) + I'(\theta)(\theta_M-\theta) + O(M^{-2}),\]

with \(I'(\theta) = \log\frac{\theta(1-p)}{p(1-\theta)} = \lambda^*\).
Hence

\[e^{-M\cdot I(\theta_M)} = e^{-M\cdot I(\theta)} \cdot \exp\!\big(-M\lambda^*(\theta_M-\theta) + O(1/M)\big).\]

Since \(\theta_M - \theta = O(1/M)\), the term
\(-M\lambda^*(\theta_M-\theta) = O(1)\). Combining all \(O(1/M)\)
errors:

\[
a_k = \frac{\exp\!\big(-M\cdot I(\theta)\big)}{\sqrt{2\pi M\cdot\theta(1-\theta)}}
\cdot \big(1 + \varepsilon_M^{(1)}\big), \qquad |\varepsilon_M^{(1)}| \leq \frac{C_2}{M}. \tag{A.1}
\]

\emph{Step 3: Ratio of consecutive terms and geometric decay.}

For \(j \geq k\),

\[\frac{a_{j+1}}{a_j} = \frac{M-j}{j+1}\cdot\frac{p}{1-p}.\]

For the term at \(j = k\):

\[\frac{a_{k+1}}{a_k} = \frac{M-k}{k+1}\cdot\frac{p}{1-p}
= e^{-\lambda_k^*}\cdot\frac{k}{k+1}, \qquad
\lambda_k^* = \log\frac{k(1-p)}{p(M-k)}.\]

Since \(k = M\theta_M\), we have \(\lambda_k^* = \lambda^* + O(1/M)\),
and \(k/(k+1) = 1 - 1/k = 1 + O(1/M)\). Thus

\[\frac{a_{k+1}}{a_k} = e^{-\lambda^*}\big(1 + O(1/M)\big).\]

Moreover, for all \(j \geq k\), the ratio is decreasing in \(j\) (since
\((M-j)/(j+1)\) decreases with \(j\)), so

\[\frac{a_{j+1}}{a_j} \leq \frac{a_{k+1}}{a_k} \leq e^{-\lambda^*}\big(1 + \rho_M\big),\]

where \(\rho_M \to 0\) as \(M \to \infty\) and \(|\rho_M| \leq C_3/M\)
for an explicit \(C_3\).

\emph{Step 4: Summing the tail.}

We bound \(P_M(\theta) = a_k + a_{k+1} + \cdots + a_M\) from both sides.

\textbf{Upper bound}:

\[
P_M(\theta) \leq a_k \sum_{j=0}^\infty \big(e^{-\lambda^*}(1+\rho_M)\big)^j
= \frac{a_k}{1 - e^{-\lambda^*}(1+\rho_M)}.
\]

\textbf{Lower bound}:

\[
P_M(\theta) \geq a_k \sum_{j=0}^{J} \big(e^{-\lambda^*}(1-\tilde\rho_M)\big)^j
\]

for any finite \(J\) (where \(\tilde\rho_M\) accounts for the lower
bound on the ratio). Taking \(J \to \infty\) (valid since the series
converges), we have

\[
P_M(\theta) \geq \frac{a_k}{1 - e^{-\lambda^*}(1-\tilde\rho_M)}.
\]

Combining both bounds and using \(|\rho_M|, |\tilde\rho_M| \leq C_3/M\):

\[
P_M(\theta) = \frac{a_k}{1 - e^{-\lambda^*}} \cdot \big(1 + \varepsilon_M^{(2)}\big),
\qquad |\varepsilon_M^{(2)}| \leq \frac{C_4}{M}. \tag{A.2}
\]

\emph{Step 5: Final assembly.}

Substituting (A.1) into (A.2):

\[
P_M(\theta) = \frac{\exp\!\big(-M\cdot I(\theta)\big)}
{(1-e^{-\lambda^*})\,\sqrt{2\pi M\cdot\theta(1-\theta)}}
\cdot \big(1 + \varepsilon_M\big), \qquad |\varepsilon_M| \leq \frac{C(p,\theta)}{M}.
\]

\textbf{Lattice correction factor.} The factor
\((1-e^{-\lambda^*})^{-1}\) is the \textbf{lattice correction}
characteristic of lattice distributions (span \(h=1\)). For large
\(\lambda^*\) (large deviations), \(e^{-\lambda^*}\) is small and
\((1-e^{-\lambda^*})^{-1} \approx 1 + e^{-\lambda^*}\) differs
negligibly from \(1\). For moderate \(\lambda^*\), the correction
matters for exact constants.

\textbf{Relation to the \(1/\lambda^*\) form.} Since
\(1-e^{-\lambda^*} = \lambda^* - \lambda^{*2}/2 + \lambda^{*3}/6 - \cdots\),

\[\frac{1}{1-e^{-\lambda^*}} = \frac{1}{\lambda^*}\left(1 + \frac{\lambda^*}{2} + \frac{\lambda^{*2}}{12} - \frac{\lambda^{*4}}{720} + \cdots\right).\]

Both forms are asymptotically equivalent as \(\lambda^* \to 0\) (i.e.,
as \(\theta \downarrow p\)), but differ by a constant factor for fixed
\(\lambda^* > 0\). For consistency with the proof architecture (Section
2.4 of the parent document), we write

\[
\boxed{\;
\mathbb{P}(\bar{X}_M \geq \theta)
= \frac{\exp(-M \cdot \text{KL}(\theta \| p))}
{\lambda^* \sqrt{2\pi M \cdot \theta(1-\theta)}}
\cdot \frac{\lambda^*}{1-e^{-\lambda^*}}
\cdot \big(1 + O(1/M)\big).\;}
\]

The factor \(\lambda^*/(1-e^{-\lambda^*})\) is a known constant in
\((0, \infty)\) and is absorbed into the leading constant in the F1
expansion (where it appears in both numerator and denominator of ratios,
so its effect can be tracked explicitly if needed). \(\square\)

\textbf{Remark (Connection to the Bahadur-Rao theorem).} The general
Bahadur-Rao (1960) theorem for non-lattice distributions yields the
\(1/\lambda^*\) form directly; the lattice correction
\((1-e^{-\lambda^*})^{-1}\) appears for lattice distributions like
Bernoulli. Both forms agree in the \(O(1/M)\) asymptotic expansion up to
a constant factor, and the \(O(1/M)\) error bound above holds in either
case.

\subsubsection{\texorpdfstring{A.4.3 Explicit Constant for the
\(O(1/M)\)
Bound}{A.4.3 Explicit Constant for the O(1/M) Bound}}\label{a.4.3-explicit-constant-for-the-o1m-bound}

\textbf{Proposition A.5 (Explicit \(O(1/M)\) bound for Bernoulli)}. For
\(p \in (0,1)\), \(\theta > p\), let
\(\lambda^* = \log\frac{\theta(1-p)}{p(1-\theta)}\),
\(\sigma^2 = \theta(1-\theta)\), and \(k = \lceil M\theta\rceil\). Then
for \(M\) large enough that \(\min(k, M-k) \geq 1\),

\[
\Bigg|
\frac{\mathbb{P}(\bar{X}_M \geq \theta) \cdot (1-e^{-\lambda^*}) \sqrt{2\pi M}\, \sigma}
{e^{-M\cdot\text{KL}(\theta\|p)}} - 1
\Bigg| \leq \frac{C_{\text{Stirling}}(p,\theta)}{M},
\]

where

\[C_{\text{Stirling}}(p,\theta) = \frac{1}{4\min(\theta, 1-\theta)}
+ \frac{C_1(p,\theta)}{\sqrt{2\pi}} + \frac{C_3(p,\theta)}{(1-e^{-\lambda^*})^2},\]

and \(C_1, C_3\) are the explicit constants from Steps 2 and 4 of the
proof above, bounded in terms of \(p\) and \(\theta\) only.

\emph{Proof.} The Stirling remainder contributes
\(|\delta_M^{(1)}| \leq \frac{1}{4\min(k,M-k)}\), the \(k/M \to \theta\)
replacement contributes \(O(1/M)\) by Lipschitz continuity of \(I\) and
the square-root factor, and the geometric-series approximation
contributes \(|\varepsilon_M^{(2)}|
\leq C_3/(M(1-e^{-\lambda^*}))\). Collecting these gives the bound.
\(\square\)

This constant is fully explicit: for any given \(p\) and \(\theta\), all
quantities can be evaluated numerically.

\begin{center}\rule{0.5\linewidth}{0.5pt}\end{center}

\subsection{\texorpdfstring{A.5 Lower Tail Formula
(\(\theta < p\))}{A.5 Lower Tail Formula (\textbackslash theta \textless{} p)}}\label{a.5-lower-tail-formula-theta-p}

For the \textbf{lower tail} \(\mathbb{P}(\bar{X}_M \leq \theta)\) with
\(\theta < p\), we use the symmetry transformation
\(Y_i = 1 - X_i \sim \text{Bern}(1-p)\). Then

\[\bar{X}_M \leq \theta \;\Longleftrightarrow\; \frac{1}{M}\sum_{i=1}^M Y_i \geq 1-\theta,\]

and \(1-\theta > 1-p\) since \(\theta < p\). Applying Theorem A.4 to the
\(Y_i\) sequence gives the following.

\textbf{Theorem A.6 (Lower tail Bahadur-Rao for Bernoulli)}. Let
\(X_1,\ldots,X_M \stackrel{i.i.d.}{\sim} \text{Bernoulli}(p)\) and fix
\(\theta < p\). Then

\[
\mathbb{P}(\bar{X}_M \leq \theta) = \frac{\exp(-M \cdot \text{KL}(\theta \| p))}
{|\lambda^*| \sqrt{2\pi M \cdot \theta(1-\theta)}}
\Bigl(1 + \varepsilon_M'\Bigr),
\]

where

\[\lambda^* = \log\frac{\theta(1-p)}{p(1-\theta)} < 0, \qquad
|\lambda^*| = \log\frac{p(1-\theta)}{\theta(1-p)},\]

and \(|\varepsilon_M'| \leq C'(p,\theta)/M\) for an explicit constant
\(C'\).

\emph{Proof.} Apply Theorem A.4 to \(Y_i = 1-X_i \sim \text{Bern}(1-p)\)
with threshold \(1-\theta > 1-p\). The saddlepoint for the \(Y\)-problem
is

\[\lambda_Y^* = \log\frac{(1-\theta)(1-(1-p))}{(1-p)(1-(1-\theta))}
= \log\frac{(1-\theta)p}{(1-p)\theta} = -\lambda^* > 0.\]

The rate function:

\[\text{KL}(1-\theta\|1-p) = (1-\theta)\log\frac{1-\theta}{1-p} + \theta\log\frac{\theta}{p} = \text{KL}(\theta\|p).\]

The variance: \((1-\theta)\theta = \theta(1-\theta)\), unchanged.
Substituting into Theorem A.4 gives the result. \(\square\)

\textbf{Application to FNR.} In the noise-detection problem where
\(e_m \sim \text{Bern}(p_1)\) under H\(_1\) and \(p_1 > \theta\):

\[\text{FNR}_M = \mathbb{P}(C_M \leq \theta \mid \text{H}_1) \sim
\frac{\exp(-M \cdot \text{KL}(\theta \| p_1))}{|\lambda_1^*| \sqrt{2\pi M \cdot \theta(1-\theta)}},\]

where \(\lambda_1^* = \log\frac{\theta(1-p_1)}{p_1(1-\theta)} < 0\),
\(|\lambda_1^*| = \log\frac{p_1(1-\theta)}{\theta(1-p_1)}\).

\begin{center}\rule{0.5\linewidth}{0.5pt}\end{center}

\subsection{A.6 Non-i.i.d. Extension}\label{a.6-non-i.i.d.-extension}

\subsubsection{A.6.1 Setting}\label{a.6.1-setting}

Consider independent but \textbf{not identically distributed} Bernoulli
variables:

\[X_m \sim \text{Bernoulli}(p_m), \qquad m = 1,\ldots,M,\]

with \(p_m \in [p_{\min}, p_{\max}] \subset (0,1)\) bounded away from 0
and 1. Define the weighted sample mean
\(C_M = \frac{1}{M}\sum_{m=1}^M X_m\) and a threshold
\(\theta > \limsup_{M\to\infty} \frac{1}{M}\sum p_m\).

\subsubsection{A.6.2 What Changes}\label{a.6.2-what-changes}

\begin{enumerate}
\def\labelenumi{\arabic{enumi}.}
\item
  \textbf{Cumulant generating function} (now depends on \(M\)):

  \[\psi_M(\lambda) = \frac{1}{M}\sum_{m=1}^M \log(1-p_m + p_m e^\lambda).\]
\item
  \textbf{Saddlepoint} \(\lambda_M^*\) solves
  \(\psi_M'(\lambda) = \theta\):

  \[\frac{1}{M}\sum_{m=1}^M \frac{p_m e^{\lambda_M^*}}{1-p_m + p_m e^{\lambda_M^*}} = \theta.\]

  This is no longer a closed form, but \(\lambda_M^*\) is bounded and
  strictly positive.
\item
  \textbf{Rate function} (Gartner-Ellis):

  \[I_M(\theta) = \sup_\lambda\{\lambda\theta - \psi_M(\lambda)\}
  = \lambda_M^*\theta - \psi_M(\lambda_M^*).\]

  This is \textbf{not} equal to \(\text{KL}(\theta\|p)\) for any single
  \(p\).
\item
  \textbf{Tilted distribution}: Under \(\mathbb{P}_{\lambda_M^*}\), the
  \(X_m\) remain independent with

  \[\mathbb{P}_{\lambda_M^*}(X_m=1) = \frac{p_m e^{\lambda_M^*}}{1-p_m + p_m e^{\lambda_M^*}} =: q_m(\theta).\]

  The tilted mean is \(\frac{1}{M}\sum q_m(\theta) = \theta\) by
  construction. The tilted variance is

  \[\sigma_M^2(\theta) = \psi_M''(\lambda_M^*) = \frac{1}{M}\sum_{m=1}^M
  \frac{p_m(1-p_m)e^{\lambda_M^*}}{(1-p_m + p_m e^{\lambda_M^*})^2}
  = \frac{1}{M}\sum_{m=1}^M q_m(\theta)(1-q_m(\theta)).\]

  Note: \(q_m(\theta) \in (p_{\min}', p_{\max}')\) bounded away from
  0,1.
\item
  \textbf{Bahadur-Rao analogue} (heterogeneous case): Under the
  Lindeberg-Feller CLT for the tilted measure,

  \[\mathbb{P}(C_M \geq \theta) = \frac{\exp(-M \cdot I_M(\theta))}
  {\lambda_M^* \sqrt{2\pi M \cdot \sigma_M^2(\theta)}}
  \Bigl(1 + O(1/M)\Bigr),\]

  provided the \(p_m\) sequence is such that the \(q_m(\theta)\) satisfy
  a Lyapunov or Lindeberg condition (which holds automatically when
  \(p_m \in [p_{\min}, p_{\max}]\)).

  The \(O(1/M)\) constant depends on \(\max_m p_m(1-p_m)\) and
  \(\min_m p_m(1-p_m)\).
\end{enumerate}

\subsubsection{A.6.3 Summary of Changes}\label{a.6.3-summary-of-changes}

\begin{longtable}[]{@{}
  >{\raggedright\arraybackslash}p{(\linewidth - 4\tabcolsep) * \real{0.2500}}
  >{\raggedright\arraybackslash}p{(\linewidth - 4\tabcolsep) * \real{0.3250}}
  >{\raggedright\arraybackslash}p{(\linewidth - 4\tabcolsep) * \real{0.4250}}@{}}
\toprule\noalign{}
\begin{minipage}[b]{\linewidth}\raggedright
Quantity
\end{minipage} & \begin{minipage}[b]{\linewidth}\raggedright
i.i.d. case
\end{minipage} & \begin{minipage}[b]{\linewidth}\raggedright
Non-i.i.d. case
\end{minipage} \\
\midrule\noalign{}
\endhead
\bottomrule\noalign{}
\endlastfoot
Saddlepoint \(\lambda^*\) & Closed form
\(\log\frac{\theta(1-p)}{p(1-\theta)}\) & Solves
\(\frac1M\sum \frac{p_m e^\lambda}{1-p_m+p_m e^\lambda} = \theta\) \\
Rate function \(I(\theta)\) & \(\text{KL}(\theta\|p)\) &
\(\lambda_M^*\theta - \frac1M\sum\log(1-p_m+p_m e^{\lambda_M^*})\) \\
Tilted variance \(\sigma^2(\theta)\) & \(\theta(1-\theta)\) &
\(\frac1M\sum q_m(\theta)(1-q_m(\theta))\) \\
Leading constant denominator & \(\lambda^*\) & \(\lambda_M^*\) \\
\end{longtable}

The \textbf{functional form} of the Bahadur-Rao expansion remains
identical; only the constants \((\lambda^*, \sigma^2, I)\) are replaced
by their \(M\)-dependent counterparts.

\begin{center}\rule{0.5\linewidth}{0.5pt}\end{center}

\subsection{A.7 Summary for Lemma A}\label{a.7-summary-for-lemma-a}

\begin{longtable}[]{@{}
  >{\raggedright\arraybackslash}p{(\linewidth - 2\tabcolsep) * \real{0.1887}}
  >{\raggedright\arraybackslash}p{(\linewidth - 2\tabcolsep) * \real{0.8113}}@{}}
\toprule\noalign{}
\begin{minipage}[b]{\linewidth}\raggedright
Quantity
\end{minipage} & \begin{minipage}[b]{\linewidth}\raggedright
Formula for i.i.d. Bern(\(p\)), \(\theta > p\)
\end{minipage} \\
\midrule\noalign{}
\endhead
\bottomrule\noalign{}
\endlastfoot
CGF \(\psi(\lambda)\) & \(\log(1-p+pe^\lambda)\) \\
Saddlepoint \(\lambda^*\) &
\(\log\frac{\theta(1-p)}{p(1-\theta)} > 0\) \\
Rate function \(I(\theta)\) &
\(\text{KL}(\theta\|p) = \theta\log\frac{\theta}{p} + (1-\theta)\log\frac{1-\theta}{1-p}\) \\
Tilted mean & \(\theta\) \\
Tilted variance \(\sigma^2(\theta)\) & \(\theta(1-\theta)\) \\
Upper tail (\(\theta > p\)) &
\(\displaystyle \frac{e^{-M\cdot\text{KL}(\theta\|p)}} {\lambda^*\sqrt{2\pi M\theta(1-\theta)}} (1+O(1/M))\) \\
Lower tail (\(\theta < p\)) &
\(\displaystyle \frac{e^{-M\cdot\text{KL}(\theta\|p)}} {|\lambda^*|\sqrt{2\pi M\theta(1-\theta)}} (1+O(1/M))\) \\
\(O(1/M)\) constant & Explicit in terms of \(p,\theta\) (Proposition
A.5) \\
\end{longtable}

\begin{center}\rule{0.5\linewidth}{0.5pt}\end{center}

\section{Lemma B: F1 Asymptotic Expansion with Higher-Order
Terms}\label{lemma-b-f1-asymptotic-expansion-with-higher-order-terms}

\begin{quote}
\textbf{Target}: Starting from the SCX F1 formula, derive the exact
asymptotic expansion of \(1-\text{F1}\) in terms of FPR and FNR, verify
the expansion's validity at the boundaries \(\eta \to 0\) and
\(\eta \to 1/2\), and substitute the Bahadur-Rao asymptotics to obtain
the leading constant.
\end{quote}

\begin{center}\rule{0.5\linewidth}{0.5pt}\end{center}

\subsection{B.1 Setup and Notation}\label{b.1-setup-and-notation}

Recall the SCX noise-detection setup:

\begin{itemize}
\tightlist
\item
  \(\eta = \mathbb{P}(\text{noise})\): prior probability of a noisy
  sample.
\item
  \(\text{TPR} = 1 - \text{FNR}_M\): true positive rate (detection
  power).
\item
  \(\text{FPR} = \text{FPR}_M\): false positive rate (type I error).
\item
  \(C_M = \frac{1}{M}\sum e_m\): consensus score.
\end{itemize}

The F1 score is defined as the harmonic mean of precision and recall:

\[\text{F1} = \frac{2\eta \cdot \text{TPR}}{\eta(1+\text{TPR}) + (1-\eta)\text{FPR}}.\]

\emph{Verification of the formula.} Using the standard definitions:

\[\text{Precision} = \frac{\eta \cdot \text{TPR}}{\eta \cdot \text{TPR} + (1-\eta)\text{FPR}}, \qquad
\text{Recall} = \text{TPR},\]

we compute

\[
\begin{aligned}
\text{F1} &= \frac{2 \cdot \text{Precision} \cdot \text{Recall}}{\text{Precision} + \text{Recall}} \\
&= \frac{2 \cdot \frac{\eta\cdot\text{TPR}}{\eta\cdot\text{TPR} + (1-\eta)\text{FPR}} \cdot \text{TPR}}
{\frac{\eta\cdot\text{TPR}}{\eta\cdot\text{TPR} + (1-\eta)\text{FPR}} + \text{TPR}} \\
&= \frac{2\eta\cdot\text{TPR}^2}{\eta\cdot\text{TPR} + (1-\eta)\text{FPR}}
\cdot \frac{1}{\frac{\eta\cdot\text{TPR}}{\eta\cdot\text{TPR} + (1-\eta)\text{FPR}} + \text{TPR}} \\
&= \frac{2\eta\cdot\text{TPR}^2}{\eta\cdot\text{TPR} + \text{TPR}(\eta\cdot\text{TPR} + (1-\eta)\text{FPR})} \\
&= \frac{2\eta\cdot\text{TPR}}{\eta + \eta\cdot\text{TPR} + (1-\eta)\text{FPR}} \\
&= \frac{2\eta\cdot\text{TPR}}{\eta(1+\text{TPR}) + (1-\eta)\text{FPR}}.
\end{aligned}
\]

This matches the expression in the proof architecture. No algebraic
error is present -- the denominator is always positive for \(\eta > 0\).

\subsection{\texorpdfstring{B.2 Exact Expression for
\(1 - \text{F1}\)}{B.2 Exact Expression for 1 - \textbackslash text\{F1\}}}\label{b.2-exact-expression-for-1---textf1}

Substituting \(\text{TPR} = 1 - \text{FNR}\):

\[
\begin{aligned}
\text{F1} &= \frac{2\eta(1-\text{FNR})}
{\eta(2-\text{FNR}) + (1-\eta)\text{FPR}}.
\end{aligned}
\]

Thus

\[
\begin{aligned}
1 - \text{F1} &= 1 - \frac{2\eta(1-\text{FNR})}
{\eta(2-\text{FNR}) + (1-\eta)\text{FPR}} \\
&= \frac{\eta(2-\text{FNR}) + (1-\eta)\text{FPR} - 2\eta(1-\text{FNR})}
{\eta(2-\text{FNR}) + (1-\eta)\text{FPR}} \\
&= \frac{\eta\cdot\text{FNR} + (1-\eta)\text{FPR}}
{\eta(2-\text{FNR}) + (1-\eta)\text{FPR}}.
\end{aligned}
\]

\textbf{Sanity checks:} -
\(\text{FNR} = \text{FPR} = 0 \implies 1-\text{F1} = 0\). -
\(\eta = 0\): \(1-\text{F1} = \text{FPR}/\text{FPR} = 1\) (F1 = 0, no
positives possible). - \(\eta = 1\):
\(1-\text{F1} = \text{FNR}/(2-\text{FNR}) \approx \text{FNR}/2\). -
Denominator
\(= 2\eta - \eta\cdot\text{FNR} + (1-\eta)\text{FPR} \geq \eta > 0\)
(strictly positive for \(\eta > 0\)).

\begin{center}\rule{0.5\linewidth}{0.5pt}\end{center}

\subsection{B.3 First-Order Expansion}\label{b.3-first-order-expansion}

Let

\[\alpha = \eta\cdot\text{FNR} + (1-\eta)\text{FPR}, \qquad
\beta = 2\eta - \eta\cdot\text{FNR} + (1-\eta)\text{FPR}.\]

Then

\[1 - \text{F1} = \frac{\alpha}{\beta}.\]

Factor \(\beta = 2\eta(1 - \gamma)\) where

\[\gamma := \frac{\eta\cdot\text{FNR} - (1-\eta)\text{FPR}}{2\eta}
= \frac{\text{FNR}}{2} - \frac{(1-\eta)\text{FPR}}{2\eta}.\]

Thus

\[1 - \text{F1} = \frac{\alpha}{2\eta} \cdot \frac{1}{1 - \gamma}.\]

For \(|\gamma| < 1\) (which holds for all sufficiently large \(M\) since
\(\text{FNR},\text{FPR} \to 0\)), the geometric series converges:

\[\frac{1}{1 - \gamma} = \sum_{k=0}^\infty \gamma^k.\]

Let

\[A := \frac{\text{FNR}}{2}, \qquad B := \frac{(1-\eta)\text{FPR}}{2\eta}.\]

Then \(\alpha/(2\eta) = A + B\), and \(\gamma = A - B\).

\textbf{First-order expansion:}

\[1 - \text{F1} = A + B + R,\]

where

\[R = (A+B) \cdot \frac{\gamma}{1-\gamma} = \frac{(A+B)(A-B)}{1-(A-B)}.\]

\textbf{Theorem B.1 (First-order expansion with error bound)}. For
\(|\gamma| = |A-B| < 1\),

\[1 - \text{F1} = \frac{\text{FNR}}{2} + \frac{(1-\eta)}{2\eta}\,\text{FPR} + R,\]

where the remainder satisfies

\[|R| \leq \frac{(\text{FNR}/2 + (1-\eta)\text{FPR}/(2\eta))^2}
{1 - \text{FNR}/2 - (1-\eta)\text{FPR}/(2\eta)}.\]

\emph{Proof.} We have \(|A+B| \leq r\) and \(|A-B| \leq r\) where
\(r = A+B = \text{FNR}/2 + (1-\eta)\text{FPR}/(2\eta)\). Then

\[|R| = \frac{(A+B)|A-B|}{1-|A-B|} \leq \frac{r^2}{1-r}.\]

For sufficiently large \(M\), \(r \leq 1/2\) and \(|R| \leq 2r^2\). This
gives the bound. \(\square\)

\textbf{Corollary B.2 (Quadratic bound)}. For \(M\) large enough that
\(r \leq 1/2\),

\[|R| \leq 2\left(\frac{\text{FNR}}{2} + \frac{(1-\eta)\text{FPR}}{2\eta}\right)^2
\leq \frac{\text{FNR}^2}{2} + \frac{(1-\eta)\text{FNR}\cdot\text{FPR}}{\eta}
+ \frac{(1-\eta)^2\text{FPR}^2}{2\eta^2}.\]

In particular,
\(R = O(\max(\text{FPR}^2, \text{FNR}^2, \text{FPR}\cdot\text{FNR}))\).

\begin{center}\rule{0.5\linewidth}{0.5pt}\end{center}

\subsection{B.4 Second-Order Expansion (Explicit
Terms)}\label{b.4-second-order-expansion-explicit-terms}

Expanding to second order in \(\gamma\):

\[
\begin{aligned}
1 - \text{F1} &= (A+B)(1 + \gamma + \gamma^2 + \gamma^3 + \cdots) \\
&= (A+B) + (A+B)(A-B) + (A+B)(A-B)^2 + O(r^4).
\end{aligned}\]

\textbf{Explicit second-order terms:}

\[
\begin{aligned}
1 - \text{F1} &= \frac{\text{FNR}}{2} + \frac{(1-\eta)\text{FPR}}{2\eta} \\
&\quad + \frac{\text{FNR}^2}{4} - \frac{(1-\eta)^2\text{FPR}^2}{4\eta^2} \\
&\quad + \big(\tfrac{\text{FNR}}{2} + \tfrac{(1-\eta)\text{FPR}}{2\eta}\big)
\big(\tfrac{\text{FNR}}{2} - \tfrac{(1-\eta)\text{FPR}}{2\eta}\big)^2 \\
&\quad + O(r^4).
\end{aligned}
\]

\textbf{Explicit \(O(\cdot)\) constant in terms of
\(\max(\text{FPR}^2, \text{FNR}^2, \text{FPR}\cdot\text{FNR})\)}:

From Corollary B.2, when \(r \leq 1/2\):

\[|R| \leq \frac{\text{FNR}^2}{2} + \frac{(1-\eta)\text{FNR}\cdot\text{FPR}}{\eta}
+ \frac{(1-\eta)^2\text{FPR}^2}{2\eta^2}.\]

Let
\(M_2 = \max(\text{FNR}^2, \text{FPR}^2, \text{FNR}\cdot\text{FPR})\).
Then

\[|R| \leq \left(\frac{1}{2} + \frac{1-\eta}{\eta} + \frac{(1-\eta)^2}{2\eta^2}\right) M_2
= \frac{1}{2\eta^2}\big(\eta^2 + 2\eta(1-\eta) + (1-\eta)^2\big) M_2
= \frac{1}{2\eta^2} \cdot 1 \cdot M_2,\]

where we used
\(\eta^2 + 2\eta(1-\eta) + (1-\eta)^2 = (\eta + (1-\eta))^2 = 1\).

Thus

\[|R| \leq \frac{1}{2\eta^2} \cdot \max(\text{FNR}^2, \text{FPR}^2, \text{FNR}\cdot\text{FPR}).\]

This gives the explicit constant \(C = 1/(2\eta^2)\) for the
\(O(\cdot)\) remainder.

\begin{center}\rule{0.5\linewidth}{0.5pt}\end{center}

\subsection{B.5 Substituting the Bahadur-Rao
Asymptotics}\label{b.5-substituting-the-bahadur-rao-asymptotics}

Now substitute the leading-order Bahadur-Rao expansions:

\[\begin{aligned}
\text{FPR}_M &\sim \frac{e^{-M\kappa_0}}{\lambda_0^* \sqrt{2\pi M \theta(1-\theta)}}, \quad
\kappa_0 = \text{KL}(\theta\|p_0),\; \lambda_0^* = \log\frac{\theta(1-p_0)}{p_0(1-\theta)} > 0, \\
\text{FNR}_M &\sim \frac{e^{-M\kappa_1}}{|\lambda_1^*| \sqrt{2\pi M \theta(1-\theta)}}, \quad
\kappa_1 = \text{KL}(\theta\|p_1),\; |\lambda_1^*| = \log\frac{p_1(1-\theta)}{\theta(1-p_1)} > 0.
\end{aligned}\]

\textbf{Theorem B.3 (Leading term of \(1-\text{F1}\) with Bahadur-Rao)}.

\[
1 - \text{F1} = \frac{1}{\sqrt{2\pi M \cdot \theta(1-\theta)}}
\left[\frac{e^{-M\kappa_1}}{2\,|\lambda_1^*|} + \frac{1-\eta}{2\eta}\cdot\frac{e^{-M\kappa_0}}{\lambda_0^*}\right]
+ O\!\left(\frac{e^{-2M\min(\kappa_0,\kappa_1)}}{M}\right).
\]

\emph{Proof.} Substitute the leading terms of FPR and FNR into the
first-order expansion from Theorem B.1. The remainder \(R\) from
Corollary B.2 is bounded by
\(O(\max(\text{FPR}^2,\text{FNR}^2)) = O(e^{-2M\min(\kappa_0,\kappa_1)}/M)\),
since each of FPR\(^2\) and FNR\(^2\) decays with twice the exponent and
an extra \(1/M\) factor from the square. \(\square\)

\subsubsection{B.5.1 Optimal Threshold: Chernoff
Point}\label{b.5.1-optimal-threshold-chernoff-point}

The optimal threshold \(\theta^*\) balances the two exponents:

\[\kappa = \text{KL}(\theta^*\|p_0) = \text{KL}(\theta^*\|p_1) = C(\text{Bern}(p_0), \text{Bern}(p_1)),\]

where \(C(\cdot,\cdot)\) is the Chernoff information. At
\(\theta = \theta^*\):

\[1 - \text{F1} \sim \frac{e^{-M\kappa}}{\sqrt{2\pi M \cdot \theta^*(1-\theta^*)}} \cdot
\left[\frac{1}{2|\lambda_1^*(\theta^*)|} + \frac{1-\eta}{2\eta \cdot \lambda_0^*(\theta^*)}\right].\]

\begin{center}\rule{0.5\linewidth}{0.5pt}\end{center}

\subsection{B.6 Verification of the Expansion at Special
Cases}\label{b.6-verification-of-the-expansion-at-special-cases}

\subsubsection{\texorpdfstring{B.6.1 Case \(\eta \to 0\) (Sparse
Noise)}{B.6.1 Case \textbackslash eta \textbackslash to 0 (Sparse Noise)}}\label{b.6.1-case-eta-to-0-sparse-noise}

As \(\eta \to 0\), the coefficient \((1-\eta)/(2\eta) \to \infty\), so
the FPR term dominates both the FNR term and the expansion itself. The
convergence condition \(|\gamma| < 1\) becomes:

\[|\gamma| = \left|\frac{\text{FNR}}{2} - \frac{(1-\eta)\text{FPR}}{2\eta}\right| < 1.\]

For small \(\eta\), this is dominated by the second term, requiring
\(\text{FPR} < 2\eta/(1-\eta) \approx 2\eta\).

Using the Bahadur-Rao approximation
\(\text{FPR} \approx e^{-M\kappa_0}/(\lambda_0^*\sqrt{M})\), the
condition becomes:

\[\frac{e^{-M\kappa_0}}{\lambda_0^*\sqrt{2\pi M \theta(1-\theta)}} \lesssim 2\eta.\]

\textbf{Definition (Threshold \(\eta_{\min}\))}. Define

\[\eta_{\min}(M) = \frac{1}{2} \cdot \frac{e^{-M\kappa_0}}{\lambda_0^*\sqrt{2\pi M \theta(1-\theta)}}.\]

For \(\eta \ll \eta_{\min}(M)\), the geometric series \(\sum \gamma^k\)
converges slowly (or diverges), and the polynomial expansion in FPR/FNR
is no longer valid. In this regime, a different asymptotic expression is
needed:

\[\begin{aligned}
1 - \text{F1} &= \frac{\eta\cdot\text{FNR} + (1-\eta)\text{FPR}}
{\eta(2-\text{FNR}) + (1-\eta)\text{FPR}} \\
&\approx \frac{\text{FPR}}{\text{FPR} + \eta/(1-\eta)} \qquad (\text{neglecting FNR terms}) \\
&= \frac{1}{1 + \eta/[(1-\eta)\text{FPR}]}.
\end{aligned}\]

When \(\eta \ll \eta_{\min}\), we have \(\eta \ll \text{FPR}\), so
\(1-\text{F1} \approx 1\) (F1 \(\approx\) 0). This is intuitive: with
extremely rare noise, even a tiny FPR destroys precision.

\textbf{Practical implication:} The expansion
\(1-\text{F1} \approx \text{FNR}/2 + (1-\eta)\text{FPR}/(2\eta)\) is
valid whenever

\[\eta \gtrsim \eta_{\min}(M) = \frac{e^{-M\kappa_0}}{2\lambda_0^*\sqrt{2\pi M \theta(1-\theta)}}.\]

For realistic \(M\), this threshold is exponentially small, so the
expansion holds in almost all practical settings.

\subsubsection{\texorpdfstring{B.6.2 Case \(\eta \to 1/2\) (Balanced
Noise)}{B.6.2 Case \textbackslash eta \textbackslash to 1/2 (Balanced Noise)}}\label{b.6.2-case-eta-to-12-balanced-noise}

When \(\eta = 1/2\), the expansion simplifies dramatically:

\[\begin{aligned}
1 - \text{F1} &= \frac{\frac{1}{2}\text{FNR} + \frac{1}{2}\text{FPR}}
{\frac{1}{2}(2-\text{FNR}) + \frac{1}{2}\text{FPR}}
= \frac{\text{FNR} + \text{FPR}}{2 - \text{FNR} + \text{FPR}} \\
&= \frac{\text{FNR}+\text{FPR}}{2} \cdot \frac{1}{1 - (\text{FNR}-\text{FPR})/2}.
\end{aligned}\]

\textbf{First-order:}

\[1 - \text{F1} = \frac{\text{FNR} + \text{FPR}}{2} + O\big(\max(\text{FNR}^2,\text{FPR}^2)\big).\]

\textbf{With Bahadur-Rao substitution at \(\theta^*\):}

\[
1 - \text{F1} \sim \frac{e^{-M\kappa}}{2\sqrt{2\pi M \cdot \theta^*(1-\theta^*)}}
\left(\frac{1}{|\lambda_1^*|} + \frac{1}{\lambda_0^*}\right).
\]

In the symmetric case \(p_0 = 1-p_1\) (which implies \(\theta^* = 1/2\),
\(\lambda_0^* = |\lambda_1^*|\)), this further simplifies to:

\[1 - \text{F1} \sim \frac{e^{-M\kappa}}{\lambda_0^* \sqrt{2\pi M \cdot \theta^*(1-\theta^*)}}.\]

\subsubsection{B.6.3 Summary of Regimes}\label{b.6.3-summary-of-regimes}

\begin{longtable}[]{@{}
  >{\raggedright\arraybackslash}p{(\linewidth - 2\tabcolsep) * \real{0.2051}}
  >{\raggedright\arraybackslash}p{(\linewidth - 2\tabcolsep) * \real{0.7949}}@{}}
\toprule\noalign{}
\begin{minipage}[b]{\linewidth}\raggedright
Regime
\end{minipage} & \begin{minipage}[b]{\linewidth}\raggedright
Leading term of \(1-\text{F1}\)
\end{minipage} \\
\midrule\noalign{}
\endhead
\bottomrule\noalign{}
\endlastfoot
\(\eta \gg \eta_{\min}\) (general) &
\(\frac{\text{FNR}}{2} + \frac{1-\eta}{2\eta}\text{FPR}\) \\
\(\eta = 1/2\) (balanced) & \((\text{FNR} + \text{FPR})/2\) \\
\(\eta \ll \eta_{\min}\) (ultra-sparse) & \(\approx 1\) (expansion
breaks down) \\
\(\eta \to 1\) (dominant noise) & \(\text{FNR}/2 + o(\text{FNR})\) \\
\end{longtable}

\begin{center}\rule{0.5\linewidth}{0.5pt}\end{center}

\subsection{\texorpdfstring{B.7 Explicit Constant for the \(O(\cdot)\)
Term}{B.7 Explicit Constant for the O(\textbackslash cdot) Term}}\label{b.7-explicit-constant-for-the-ocdot-term}

Collecting the results from Sections B.3-B.5, we have the complete
asymptotic:

\textbf{Theorem B.4 (Full expansion with explicit constants)}. Let
\(\text{FPR} = \text{FPR}_M\) and \(\text{FNR} = \text{FNR}_M\) be the
error rates from the threshold test with threshold \(\theta\). Assume
\(M\) is large enough that

\[r := \frac{\text{FNR}}{2} + \frac{(1-\eta)\text{FPR}}{2\eta} \leq \frac{1}{2}.\]

Then

\[
1 - \text{F1} = \frac{\text{FNR}}{2} + \frac{1-\eta}{2\eta}\text{FPR}
+ \frac{\text{FNR}^2}{4} - \frac{(1-\eta)^2\text{FPR}^2}{4\eta^2} + R_3,
\]

where
\(|R_3| \leq 2r^3 \leq \frac{1}{4\eta^3}\big(\eta\cdot\text{FNR} + (1-\eta)\text{FPR}\big)^3\).

When \(\eta \gtrsim \eta_{\min}(M)\), substituting the Bahadur-Rao
expansions yields:

\[
1 - \text{F1} = \frac{1}{\sqrt{2\pi M\theta(1-\theta)}}
\left[\frac{e^{-M\kappa_1}}{2|\lambda_1^*|} + \frac{1-\eta}{2\eta}\frac{e^{-M\kappa_0}}{\lambda_0^*}\right]
\cdot \big(1 + O(1/M)\big) + O\!\left(\frac{e^{-2M\min(\kappa_0,\kappa_1)}}{M}\right).
\]

At the Chernoff point \(\theta^*\) where
\(\kappa_0 = \kappa_1 = \kappa\):

\[
1 - \text{F1} = \frac{e^{-M\kappa}}{\sqrt{2\pi M\theta^*(1-\theta^*)}}
\left[\frac{1}{2|\lambda_1^*|} + \frac{1-\eta}{2\eta\lambda_0^*}\right]
\cdot \big(1 + O(1/M)\big).
\]

\begin{center}\rule{0.5\linewidth}{0.5pt}\end{center}

\subsection{B.8 Verification: No Denominator
Bug}\label{b.8-verification-no-denominator-bug}

The user noted a concern about the denominator going negative in v1. We
verify:

\[\text{Denominator} = \eta(2-\text{FNR}) + (1-\eta)\text{FPR}.\]

Since \(\text{FNR} \leq 1\), we have \(2 - \text{FNR} \geq 1\), so

\[\text{Denominator} \geq \eta \cdot 1 + 0 = \eta > 0 \quad (\text{for } \eta > 0).\]

For \(\eta = 0\), the F1 formula is degenerate (no positive samples),
which is handled separately. \textbf{No algebraic error exists} in the
F1 expression used in the proof architecture.

\begin{center}\rule{0.5\linewidth}{0.5pt}\end{center}

\subsection{References}\label{references}

\begin{enumerate}
\def\labelenumi{\arabic{enumi}.}
\tightlist
\item
  Bahadur, R. R., \& Rao, R. R. (1960). On deviations of the sample
  mean. \emph{Annals of Mathematical Statistics}, 31(4), 1015-1027.
\item
  Cramer, H. (1938). Sur un nouveau theoreme-limite de la theorie des
  probabilites. \emph{Actualites Scientifiques et Industrielles}, 736,
  5-23.
\item
  Dembo, A., \& Zeitouni, O. (2010). \emph{Large Deviations Techniques
  and Applications} (2nd ed.). Springer.
\item
  Jensen, J. L. (1995). \emph{Saddlepoint Approximations}. Oxford
  University Press.
\item
  Lugannani, R., \& Rice, S. (1980). Saddle point approximation for the
  distribution of the sum of independent random variables.
  \emph{Advances in Applied Probability}, 12(2), 475-490.
\item
  Shevtsova, I. G. (2011). On the absolute constants in the
  Berry-Esseen-type inequalities for identically distributed summands.
  \emph{Doklady Mathematics}, 83(3), 320-323. {[}For the optimal
  Berry-Esseen constant \(C_{\text{BE}} \leq 0.4748\).{]}
\end{enumerate}

\end{document}
